\chapter{Testing and Results}

\section{Introduction}
Rigorous experimental testing and validation are essential to verify that the cyber-physical StepCharge AI platform satisfies its functional and non-functional requirements. This chapter presents the testing methodologies, experimental test suites, empirical observations, and performance analyses conducted across all architectural tiers. Testing spans physical hardware characterization, analog signal acquisition, firmware step detection, the 101-test automated backend Vitest suite, machine learning cross-validation, and real-time dashboard latency measurements.

\section{Hardware Testing}
Hardware evaluation examined the electromechanical properties of the piezoelectric transducer array under controlled vertical compressive impacts:
\begin{itemize}
    \item \textbf{Open-Circuit Transducer Potential:} Direct oscilloscope probing across a single 35\,mm PZT-5 disc during an unbuffered footfall produced open-circuit AC voltage peaks ranging from $14.2\,\text{V}$ to $32.6\,\text{V}$.
    \item \textbf{Rectification and Forward Drop:} Probing downstream of the 1N5819 Schottky bridge confirmed full-wave rectification into positive DC pulses, with an observed forward voltage drop of $0.34\,\text{V}$ per diode branch, verifying minimal power loss.
    \item \textbf{Overvoltage Clamping Verification:} Heavy simulated impacts ($>800\,\text{N}$) generated raw spikes approaching $42\,\text{V}$. The 5.1\,V Zener diode successfully clamped the DC rail to $5.08\,\text{V}$, preventing damage to the 5.5\,V supercapacitor and the ESP32 ADC.
\end{itemize}

\section{Voltage Measurement Testing}
The precision of the $100\,\text{k}\Omega / 100\,\text{k}\Omega$ resistive voltage dividers was calibrated against a benchtop 6.5-digit digital multimeter (DMM). Table~\ref{tab:voltage_calibration} records the calibration measurements across the operational range.

\begin{table}[H]
\centering
\small
\renewcommand{\arraystretch}{1.25}
\caption{Voltage measurement calibration and ADC conversion accuracy across operational range.}
\label{tab:voltage_calibration}
\begin{tabularx}{\textwidth}{@{} C C C C C @{}}
\toprule
\rowcolor{iceblue}
\tableheader{Applied Voltage (V)} & \tableheader{Divider Out (V)} & \tableheader{ESP32 Measured (V)} & \tableheader{Absolute Error (V)} & \tableheader{Percentage Error (\%)} \\
\midrule
\rowcolor{white}
1.000 & 0.499 & 0.992 & 0.008 & 0.80\% \\
\rowcolor{lightgrey}
2.000 & 0.998 & 1.991 & 0.009 & 0.45\% \\
\rowcolor{white}
3.000 & 1.498 & 2.988 & 0.012 & 0.40\% \\
\rowcolor{lightgrey}
4.000 & 1.996 & 3.985 & 0.015 & 0.38\% \\
\rowcolor{white}
5.000 & 2.495 & 4.981 & 0.019 & 0.38\% \\
\midrule
\rowcolor{iceblue}
\multicolumn{3}{@{}l}{\textbf{\color{deepnavy}Mean Percentage Error Across Range}} & \multicolumn{2}{r}{\textbf{\color{deepnavy}0.48\% (Calibrated)}} \\
\bottomrule
\end{tabularx}
\end{table}

The mean percentage error of $0.48\%$ across the operational storage range confirms that the 16$\times$ hardware oversampling and calibrated conversion formula deliver high measurement accuracy.

\section{Footstep Detection Testing}
The dual-threshold Schmitt trigger firmware was validated through repeated human walking trials:
\begin{itemize}
    \item \textbf{Hysteresis Verification:} A single footstep produces multiple mechanical rebound oscillations during the heel-off phase. The rising threshold of $V_{\text{high}} = 0.80\,\text{V}$ initiates pulse tracking, while the falling threshold of $V_{\text{low}} = 0.55\,\text{V}$ delays pulse termination until full release.
    \item \textbf{Refractory Lockout Effectiveness:} Without refractory lockout, 100 physical footsteps registered 128 counts (a 28\% false-positive rate due to mechanical bounce). Enabling the 250\,ms refractory lockout reduced false double-triggering to zero, yielding 100\% counting accuracy across standard cadence ranges ($1.0\text{--}2.2\,\text{steps/s}$).
    \item \textbf{Duration Filtering:} Simulated high-frequency electrical transients ($<30\,\text{ms}$) and continuous static standing ($>3\,\text{s}$) were successfully rejected by the duration boundary filter ($60\,\text{ms} \le T_{\text{pulse}} \le 2000\,\text{ms}$).
\end{itemize}

\section{ESP32 Communication Testing}
Wi-Fi resilience was tested under variable radio frequency (RF) conditions:
\begin{itemize}
    \item \textbf{Telemetry Latency:} Over 1,000 consecutive transmission cycles, the mean HTTP POST round-trip latency was $184\,\text{ms}$ ($\sigma = 38\,\text{ms}$) on a 2.4\,GHz WPA2 network.
    \item \textbf{Reconnection Recovery:} When the Wi-Fi router was temporarily powered down, the ESP32 detected socket loss within 4.0 seconds, buffered incoming footstep events in a circular queue in SRAM (capacity: 100 events), and successfully re-associated and flushed the buffer upon signal restoration with zero packet loss.
\end{itemize}

\section{Backend API and Security Testing}
The backend application was subjected to an automated test suite implemented in Vitest. All \textbf{101 unit and integration tests} executed successfully (100\% pass rate) across seven distinct test suites:
\begin{enumerate}
    \item \texttt{security.test.ts} (35 tests): Verifies authentication guards (401 on unauthenticated requests, valid cookie sessions), device ownership isolation (403 on cross-user device access), device key revocation, physical input bounds (rejection of $V > 60\,\text{V}$, NaN values, or $V_{\text{avg}} > V_{\text{peak}}$), secret masking (ensuring MongoDB connection strings, JWT secrets, and Google OAuth Client Secrets are never leaked), and fail-closed handling.
    \item \texttt{auth\_flow.test.ts} (6 tests): Validates Google OAuth callback parameter handling, CSRF state verification, JWT cookie issuance, and logout cookie clearing.
    \item \texttt{api.test.ts} (13 tests): Assesses REST route health, telemetry ingest endpoints, and device metadata queries.
    \item \texttt{backend.test.ts} (11 tests): Tests MongoDB connection resiliency, graceful shutdown hooks, and service lifecycle.
    \item \texttt{phase2\_measurement.test.ts} (17 tests): Tests dual-rail voltage calculations, supercapacitor energy equations ($E = \frac{1}{2} C V^2$), and unit conversions.
    \item \texttt{phase3\_ai\_intelligence.test.ts} (13 tests): Evaluates ML payload construction, feature bounds validation, and graceful handling of ML microservice connection timeouts.
    \item \texttt{phase4\_research\_validation.test.ts} (6 tests): Validates research dataset exports, CSV serialization, and statistical aggregation pipelines.
\end{enumerate}

\section{Database Testing}
MongoDB Atlas database performance was evaluated under concurrent write stress:
\begin{itemize}
    \item \textbf{Write Throughput:} Sustained write operations of 100 telemetry frames per second yielded an average insertion latency of $4.2\,\text{ms}$.
    \item \textbf{Indexed Query Efficiency:} Querying the most recent 100 telemetry records using the compound index \texttt{\{ deviceId: 1, timestamp: -1 \}} executed in $0.85\,\text{ms}$ on a dataset containing 500,000 historical documents, demonstrating $O(\log N)$ scalability.
\end{itemize}

\section{Machine Learning Model Evaluation}
The Random Forest footstep classifier was trained and evaluated on a ground-truth dataset comprising 360 recorded footstep impact transients collected across 6 human participants (3 male, 3 female; body masses: $52\text{--}88\,\text{kg}$). Footsteps were recorded across three controlled intensity protocols:
\begin{itemize}
    \item \textbf{LIGHT:} Gentle, shuffling walking or tiptoe impact ($V_{\text{peak}} < 1.8\,\text{V}$).
    \item \textbf{NORMAL:} Standard casual walking with typical heel-strike impact ($1.8\,\text{V} \le V_{\text{peak}} \le 3.2\,\text{V}$).
    \item \textbf{HEAVY:} Forceful walking, running impact, or weighted stomping ($V_{\text{peak}} > 3.2\,\text{V}$).
\end{itemize}

\subsection{Subject-Independent vs. Random Split Evaluation}
To validate the necessity of subject-independent cross-validation, the model was evaluated under two distinct partitioning protocols:
\begin{enumerate}
    \item \textbf{Standard Random Split (Subject-Dependent):} A naive 80/20 train/test random split produced an artificially inflated test accuracy of \textbf{94.4\%}. However, feature importance analysis revealed that the model was partially learning participant-specific cadence intervals ($\Delta T_{\text{step}}$) rather than kinetic strike force.
    \item \textbf{Stratified Group $K$-Fold (Subject-Independent):} Enforcing 5-fold cross-validation grouped strictly by \texttt{participantId} (ensuring zero participant overlap between folds) yielded a realistic, generalizable accuracy of \textbf{88.6\%}.
\end{enumerate}

\subsection{Confusion Matrix and Classification Metrics}
Table~\ref{tab:confusion_matrix} displays the aggregate confusion matrix evaluated under subject-independent Group $K$-Fold cross-validation.

\begin{table}[H]
\centering
\small
\renewcommand{\arraystretch}{1.25}
\caption{Confusion matrix under subject-independent Stratified Group $K$-Fold validation.}
\label{tab:confusion_matrix}
\begin{tabular}{@{} l c c c r @{}}
\toprule
\rowcolor{iceblue}
\tableheader{Actual Class} & \multicolumn{3}{c}{\tableheader{Predicted Class}} & \tableheader{Total Support} \\
\cmidrule(lr){2-4}
\rowcolor{iceblue}
& \textbf{\color{deepnavy}LIGHT} & \textbf{\color{deepnavy}NORMAL} & \textbf{\color{deepnavy}HEAVY} & \\
\midrule
\rowcolor{white}
\textbf{\color{deepnavy}LIGHT} & \textbf{108} & 11 & 1 & 120 \\
\rowcolor{lightgrey}
\textbf{\color{deepnavy}NORMAL} & 13 & \textbf{102} & 5 & 120 \\
\rowcolor{white}
\textbf{\color{deepnavy}HEAVY} & 0 & 11 & \textbf{109} & 120 \\
\midrule
\rowcolor{iceblue}
\textbf{\color{deepnavy}Total Predictions} & 121 & 124 & 115 & \textbf{360} \\
\bottomrule
\end{tabular}
\end{table}

The resulting class-level performance metrics are summarized in Table~\ref{tab:ml_metrics}.

\begin{table}[H]
\centering
\small
\renewcommand{\arraystretch}{1.25}
\caption{Detailed Machine Learning classification metrics per gait class.}
\label{tab:ml_metrics}
\begin{tabular}{@{} l c c c r @{}}
\toprule
\rowcolor{iceblue}
\tableheader{Impact Gait Class} & \tableheader{Precision} & \tableheader{Recall} & \tableheader{F1-Score} & \tableheader{Support (Samples)} \\
\midrule
\rowcolor{white}
\texttt{LIGHT}  & 0.893 & 0.900 & 0.896 & 120 \\
\rowcolor{lightgrey}
\texttt{NORMAL} & 0.823 & 0.850 & 0.836 & 120 \\
\rowcolor{white}
\texttt{HEAVY}  & 0.948 & 0.908 & 0.928 & 120 \\
\midrule
\rowcolor{iceblue}
\textbf{\color{deepnavy}Macro Average} & \textbf{0.888} & \textbf{0.886} & \textbf{0.887} & \textbf{360} \\
\rowcolor{iceblue}
\textbf{\color{deepnavy}Overall Accuracy} & \multicolumn{4}{c}{\textbf{\color{deepnavy}88.61\% (Subject-Independent GroupKFold)}} \\
\bottomrule
\end{tabular}
\end{table}

The model demonstrates strong discrimination, particularly on \texttt{HEAVY} impacts (Precision: $94.8\%$), with minor misclassifications occurring primarily along the continuous boundary between \texttt{LIGHT} and \texttt{NORMAL} walking.

\section{Dashboard Testing}
The React 19 web dashboard was evaluated across Chromium (Chrome v122, Edge v122) and Gecko (Firefox v123) browsers. Live telemetry streaming maintained stable 60\,FPS rendering without memory leakage over a continuous 4-hour test. The accessible profile menu was verified for correct focus trapping, outside-click closing, and instantaneous dismissal upon pressing the Escape key.

\section{Integration Testing}
End-to-end integration testing confirmed the closed feedback loop:
\begin{itemize}
    \item \textbf{Telemetry Pipeline Latency:} The duration from physical footstep impact to graphical rendering on the web dashboard averaged \textbf{86\,ms} ($12\,\text{ms}$ firmware processing + $48\,\text{ms}$ HTTPS transmission + $18\,\text{ms}$ backend and Socket.IO emission + $8\,\text{ms}$ React DOM update).
    \item \textbf{Closed-Loop Load Feedback:} When a user clicked ``Toggle LED'', the backend recorded the command intent. The ESP32 polled the intent within $1500\,\text{ms}$, energized the MOSFET gate, read the GPIO pin voltage ($3.28\,\text{V}$ confirmed high), and reported \texttt{actualState = true}. The dashboard UI rendered the active state only upon receiving this confirmation, verifying closed-loop integrity.
\end{itemize}

\section{Comprehensive Test Case Table}
Table~\ref{tab:test_cases} documents 15 formal test cases spanning physical, embedded, API, and client tiers.

\begin{landscape}
\begin{table}[p]
\centering
\scriptsize
\renewcommand{\arraystretch}{1.22}
\caption{Comprehensive system test cases, empirical verification, and validation status.}
\label{tab:test_cases}
\begin{tabularx}{\linewidth}{@{} L{0.6in} L{1.5in} L{1.8in} L{2.0in} Y M{0.7in} @{}}
\toprule
\rowcolor{iceblue}
\tableheader{Test ID} & \tableheader{Test Description} & \tableheader{Input / Trigger Action} & \tableheader{Expected Specification} & \tableheader{Empirical System Result} & \tableheader{Status} \\
\midrule
\rowcolor{white}
TC-01 & Mechanical Rectification & Footstep impact on tile & AC converted to DC with $<0.40\,\text{V}$ drop & Full-wave DC, $V_F = 0.34\,\text{V}$ measured & \passbadge \\
\rowcolor{lightgrey}
TC-02 & Overvoltage Clamping & Forceful impact ($>800\,\text{N}$) & DC voltage clamped below $5.50\,\text{V}$ & Clamped at $5.08\,\text{V}$ via 1N4733A Zener & \passbadge \\
\rowcolor{white}
TC-03 & Schmitt Trigger Step Detection & Single footstep ($V_{\text{pk}} = 2.4\,\text{V}$) & Increments step count exactly by 1 & Count incremented by exactly 1 & \passbadge \\
\rowcolor{lightgrey}
TC-04 & Refractory Bounce Rejection & Mechanical bounce $120\,\text{ms}$ post-step & Ignored during $250\,\text{ms}$ lockout window & Refractory window suppressed false trigger & \passbadge \\
\rowcolor{white}
TC-05 & Noise Duration Filter & Electrical noise spike ($20\,\text{ms}$) & Rejected as transient ($T_{\text{pulse}} < 60\,\text{ms}$) & Rejected successfully, count unchanged & \passbadge \\
\rowcolor{lightgrey}
TC-06 & ADC Linear Calibration & Apply $3.00\,\text{V}$ from precision source & ESP32 reads $3.00\,\text{V} \pm 0.05\,\text{V}$ & Calibrated reading: $2.988\,\text{V}$ ($-0.4\%$ error) & \passbadge \\
\rowcolor{white}
TC-07 & Stored Energy Computation & Storage voltage at $4.00\,\text{V}$ & $E = \frac{1}{2}(0.10)(4)^2 = 0.80\,\text{J}$ & Calculated $0.800\,\text{J}$ matching formula & \passbadge \\
\rowcolor{lightgrey}
TC-08 & Device Authentication & POST telemetry with valid SHA-256 key & HTTP 200 OK with session token & HTTP 200 OK, telemetry persisted & \passbadge \\
\rowcolor{white}
TC-09 & Invalid Key Rejection & POST telemetry with invalid API key & HTTP 403 Forbidden with security log & HTTP 403 Forbidden, request dropped & \passbadge \\
\rowcolor{lightgrey}
TC-10 & Physical Bounds Check & POST telemetry with $V_{\text{avg}} > V_{\text{peak}}$ & HTTP 400 Bad Request error payload & HTTP 400 Bad Request, rejected cleanly & \passbadge \\
\rowcolor{white}
TC-11 & Google OAuth Session & Complete OAuth 2.0 redirect callback & Issue secure HttpOnly session cookie & Cookie issued securely with \texttt{SameSite=Lax} & \passbadge \\
\rowcolor{lightgrey}
TC-12 & ML Inference Microservice & POST valid feature vector to FastAPI & Return class and confidence score & Returned \texttt{NORMAL} ($92\%$ conf) in $18\,\text{ms}$ & \passbadge \\
\rowcolor{white}
TC-13 & Real-Time Socket Broadcast & Ingest telemetry frame into backend & Connected clients receive \texttt{telemetry:update} & Pushed to React client in $<45\,\text{ms}$ & \passbadge \\
\rowcolor{lightgrey}
TC-14 & Closed-Loop Load Verification & Send load toggle command on dashboard & UI state toggles only after GPIO readback & MOSFET energized, readback $3.28\,\text{V}$ & \passbadge \\
\rowcolor{white}
TC-15 & Notification Dismissal & Click ``Mark as Read'' on alert drawer & MongoDB updated; unread count set to 0 & Persisted and synced across all tabs & \passbadge \\
\bottomrule
\end{tabularx}
\end{table}
\end{landscape}

\section{Observed Results}
Figure~\ref{fig:cap_charging} presents the empirical charging trajectory of the 0.1\,F supercapacitor over a sequence of 60 consecutive pedestrian walking steps.

\begin{figure}[H]
\centering
\begin{tikzpicture}[scale=0.92, every node/.style={scale=0.92}]
    \draw [->, thick, draw=deepnavy] (0,0) -- (10.2,0) node [right, font=\scriptsize\bfseries] {Footsteps (Count)};
    \draw [->, thick, draw=deepnavy] (0,0) -- (0,5.2) node [above, font=\scriptsize\bfseries] {Storage Voltage (V)};

    % Grid lines
    \draw [help lines, color=lightgrey!70!gray, line width=0.5pt] (0,0) grid (9.5,4.5);

    % Axis labels
    \foreach \x in {0, 10, 20, 30, 40, 50, 60}
        \draw [draw=deepnavy] (\x*9.0/60,0) -- (\x*9.0/60,-0.1) node [below, font=\scriptsize] {\x};
    \foreach \y in {0, 1, 2, 3, 4}
        \draw [draw=deepnavy] (0,\y) -- (-0.1,\y) node [left, font=\scriptsize] {\y.0};

    % Empirical Voltage Curve (charging curve)
    \draw [line width=1.4pt, profblue, smooth] plot coordinates {
        (0, 0.20)
        (0.75, 0.45)
        (1.50, 0.72)
        (2.25, 0.98)
        (3.00, 1.25)
        (3.75, 1.55)
        (4.50, 1.88)
        (5.25, 2.22)
        (6.00, 2.58)
        (6.75, 2.92)
        (7.50, 3.25)
        (8.25, 3.55)
        (9.00, 3.82)
    };

    % Final Operating Point
    \node at (9.0, 3.82) [circle, fill=profblue, inner sep=2.5pt] {};
    \node at (9.0, 4.2) [font=\scriptsize\bfseries\color{deepnavy}] {$V_{\text{storage}} = 3.82\,\text{V}$ ($E = 0.730\,\text{J}$)};

    % Target guideline
    \draw [dashed, amber, line width=1.1pt] (0,3.82) -- (9.0,3.82);

\end{tikzpicture}
\caption{Empirical supercapacitor voltage charging curve over 60 consecutive footsteps.}
\label{fig:cap_charging}
\end{figure}

The terminal storage voltage rose monotonically from an initial baseline of $0.20\,\text{V}$ to $3.82\,\text{V}$ over 60 footsteps, accumulating an estimated stored electrostatic energy of:
\begin{equation}
E = \frac{1}{2} (0.10\,\text{F}) (3.82\,\text{V})^2 \approx 0.730\,\text{Joules}
\end{equation}
yielding an average energy increment of approximately $12.1\,\text{mJ}$ per footstep.

\section{Result Analysis}
Statistical evaluation was conducted using non-parametric rank testing. Because footstep impact amplitudes exhibit non-Gaussian distributions due to biomechanical gait variability, the \textbf{Kruskal-Wallis $H$-test} was performed across the three gait classes:
\begin{equation}
H = \frac{12}{N(N+1)} \sum_{j=1}^k \frac{R_j^2}{n_j} - 3(N+1)
\end{equation}
With $N = 360$, $k = 3$, and degrees of freedom $\text{df} = 2$, the test yielded $H = 268.4$ ($p < 0.001$). The statistically significant $p$-value confirms that the three footstep intensity classes occupy distinct physical energy distributions, rejecting the null hypothesis that impact voltages arise from identical continuous distributions.

\section{Limitations}
While the experimental results validate system efficacy, four primary physical and computational limitations must be acknowledged:
\begin{enumerate}
    \item \textbf{Footwear Sole Damping:} Tests conducted with thick elastomeric running shoes attenuated peak voltage by approximately $38\%$ and lengthened pulse duration compared to hard leather soles, indicating that footwear compliance introduces low-pass mechanical filtering.
    \item \textbf{Supercapacitor Self-Discharge:} Internal leakage currents ($5\text{--}15\,\mu\text{A}$) cause stored energy to decay gradually during extended quiescent intervals ($>4\,\text{hours}$).
    \item \textbf{Unmeasured Current Provenance:} The current prototype does not incorporate an active current shunt (such as an INA219). Consequently, stored energy is derived purely from terminal voltage differentials. Instantaneous current and delivered power are not measured directly.
    \item \textbf{Substrate Scaling:} Evaluation was performed on a single isolated tile rather than an interconnected multi-tile array, limiting total cumulative energy generation.
\end{enumerate}
